\documentclass[11pt]{article}
\usepackage{docstyle}
\renewcommand\sheetname{Homework \textperiodcentered{} Human Reproduction}
\begin{document}
\sheethead{Homework}{Human Reproduction \textperiodcentered{} Year 9 Science}{about 80 minutes \textperiodcentered{} A / M / E}
Do Part 1 and Part 2 on two different days before the exam (about 40 minutes each). \textbf{Print this sheet at 100\,\%}: the graph grid is true size.


\section{Part 1 \textperiodcentered{} Sex cells, the two systems and the menstrual cycle}
\Q Write the meaning of each word.\lvl{A}
\par\sub{a} gamete
\alines{1}{a sex cell: a sperm cell or an egg cell (ovum)}
\sub{b} ovulation
\alines{1}{the release of an egg from an ovary, on about day 14 of the cycle}
\sub{c} hormone
\alines{1}{a chemical messenger, made by a gland and carried in the blood to the organs it affects}

\Q Is each part found in males only, in females only, or in both sexes? Write \textbf{M}, \textbf{F} or \textbf{B}.\lvl{A}
\par\vspace{1mm}
\begin{tabular}{@{}l@{\hspace{3mm}}l@{\hspace{12mm}}l@{\hspace{3mm}}l@{\hspace{12mm}}l@{\hspace{3mm}}l@{\hspace{12mm}}l@{\hspace{3mm}}l@{}}
bladder & \blank[10mm]{B} & cervix & \blank[10mm]{F} & epididymis & \blank[10mm]{M} & urethra & \blank[10mm]{B}\\[2mm]
oviduct & \blank[10mm]{F} & prostate gland & \blank[10mm]{M} & ovary & \blank[10mm]{F} & anus & \blank[10mm]{B}\\
\end{tabular}\par\vspace{2mm}

\Q As soon as one sperm cell has entered an egg cell, the membrane of the egg changes.\lvl{M}
\par\sub{a} State what this change does.
\alines{1}{It stops any other sperm cell from entering the egg.}
\sub{b} Explain why this is important for the number of chromosomes in the zygote.
\alines{3}{One sperm brings 23 chromosomes and the egg has 23, so the zygote has 46. A second sperm would add 23 more, and the zygote would have too many chromosomes to develop normally.}

\Q Explain why a girl cannot become pregnant before puberty.\lvl{E}
\alines{4}{Before puberty no eggs are released. At puberty, hormones start the menstrual cycle, and ovulation begins. Without an egg in the oviduct there is nothing for a sperm to fertilise, so no zygote can form.}

\Q State what happens at the menopause, and the age at which it usually happens.\lvl{A}
\alines{2}{Egg release and periods stop, so the woman can no longer become pregnant. It happens at around 50 years.}

\Q A woman measures her body temperature at the same time every morning during one cycle.
\par\vspace{1mm}
\begin{tabular}{@{}l *{8}{c}@{}}
\toprule
Day of the cycle & 2 & 6 & 10 & 13 & 15 & 18 & 22 & 26\\
Temperature (\textdegree C) & 36.4 & 36.4 & 36.5 & 36.4 & 36.8 & 36.9 & 36.9 & 36.8\\
\bottomrule
\end{tabular}\par\vspace{1mm}
\sub{a} Plot the results on the grid. Label both axes with units, choose an even scale, and join the points with straight lines.\lvl{M}
\par\vspace{1mm}
\hbox to\linewidth{\hss\GraphGrid{14}{6}{%
  \foreach \d/\t in {2/36.4,6/36.4,10/36.5,13/36.4,15/36.8,18/36.9,22/36.9,26/36.8}
    \draw ({\d*0.5-0.1},{(\t-36)*5-0.1}) -- ({\d*0.5+0.1},{(\t-36)*5+0.1}) ({\d*0.5-0.1},{(\t-36)*5+0.1}) -- ({\d*0.5+0.1},{(\t-36)*5-0.1});
  \draw (1,2.0) -- (3,2.0) -- (5,2.5) -- (6.5,2.0) -- (7.5,4.0) -- (9,4.5) -- (11,4.5) -- (13,4.0);
  \foreach \d in {0,4,...,28} \node[below] at (\d*0.5,0) {\small\d};
  \foreach \t/\y in {36.0/0,36.2/1,36.4/2,36.6/3,36.8/4,37.0/5} \node[left] at (0,\y) {\small\t};
  \node[below=4mm] at (7,0) {\small day of the cycle};
  \node[rotate=90] at (-1.35,3) {\small temperature (\textdegree C)};}\hss}
\ansnote{5mm}{Scale: 1 cm = 2 days (x), 1 cm = 0.2 \textdegree C (y), starting at 36.0 \textdegree C. Points joined with straight lines.}
\sub{b} Describe the pattern in the results.\lvl{M}
\alines{3}{From day 2 to day 13 the temperature stays at about 36.4 to 36.5 \textdegree C. Between day 13 and day 15 it rises by 0.4 \textdegree C. It then stays at about 36.8 to 36.9 \textdegree C until day 26.}
\sub{c} Use the results to suggest the day of ovulation. Explain your answer.\lvl{E}
\alines{2}{About day 14, because the temperature rises just after an egg is released and then stays higher.}

\Q A student is asked to describe the menstrual cycle.\par\vspace{1mm}
\samcard{0.9\linewidth}{\flawed{The menstrual cycle lasts about 14 days. LH makes the lining of the uterus break down. If the egg is fertilised, the woman has a period.}}\par\vspace{1mm}
\textbf{Sam's answer has 3 errors.} Find each one and write the correct science.\lvl{M}
\alines{5}{1. The cycle lasts about 28 days; day 14 is ovulation (\Lref{ovul}). 2. LH triggers ovulation; the lining breaks down when the level of progesterone falls (\Lref{horm}). 3. If the egg is fertilised, the embryo implants and the lining stays thick, so there is no period (\Lref{stages}).}

\newpage
\section{Part 2 \textperiodcentered{} Fertilisation, chromosomes, pregnancy and fertility}
\Q Write the meaning of each word.\lvl{A}
\par\sub{a} zygote
\alines{1}{the cell formed at fertilisation, when a sperm nucleus fuses with an egg nucleus}
\sub{b} implantation
\alines{1}{when the embryo sinks into the thick lining of the uterus}
\sub{c} premature
\alines{1}{born before the expected time: for a human baby, before 38 weeks}

\Q A body cell of a dog has 78 chromosomes.\lvl{A / M}
\par (a) Number of chromosomes in a dog's sperm cell: \blank[16mm]{39} \qquad (b) in a fertilised dog egg: \blank[16mm]{78}
\par\vspace{1mm}\sub{c} Explain your answer to (a).
\alines{2}{A sex cell has half the number of a body cell (one set), so that the zygote has the full number when two sex cells join.}

\Q Explain why about half of all babies are boys.\lvl{M}
\alines{3}{Every egg carries an X chromosome. Half of a man's sperm carry X and half carry Y. So XX (a girl) and XY (a boy) are equally likely at each fertilisation.}

\Q Complete the table about twins.\lvl{A}
\par\vspace{1mm}
\begin{tabular}{@{}l@{\hspace{8mm}}c@{\hspace{8mm}}c@{}}
\toprule
 & \textbf{Identical twins} & \textbf{Non-identical twins}\\
\midrule
Number of eggs fertilised & \blank[22mm]{1} & \blank[22mm]{2}\\[1.5mm]
Number of sperm cells needed & \blank[22mm]{1} & \blank[22mm]{2}\\[1.5mm]
Do they have the same genes? & \blank[22mm]{yes} & \blank[22mm]{no}\\[1.5mm]
Are they always the same sex? & \blank[22mm]{yes} & \blank[22mm]{no}\\
\bottomrule
\end{tabular}\par\vspace{2mm}

\Q The placenta has many villi and a rich blood supply. Explain how each feature helps the foetus.\lvl{E}
\alines{5}{The villi give the placenta a very large surface area, so a lot of oxygen and glucose can diffuse into the foetus's blood, and carbon dioxide and urea out of it, at the same time. The rich blood supply keeps bringing oxygen and glucose and carrying wastes away, so the difference in concentration stays large and diffusion stays fast.}

\Q A pregnant woman slips and falls on her side. Explain why the foetus is usually not hurt.\lvl{M}
\alines{2}{The foetus floats in amniotic fluid inside the amniotic sac. The fluid cushions the foetus and absorbs the bump.}

\Q Name the three stages of birth in order. State what leaves the uterus in stage 2 and in stage 3.\lvl{A}
\par Stage 1 \blank[30mm]{labour} \quad Stage 2 \blank[30mm]{delivery} \quad Stage 3 \blank[30mm]{afterbirth}
\par\vspace{2.5mm} Leaves in stage 2: \blank[40mm]{the baby} \quad Leaves in stage 3: \blank[40mm]{the placenta}\par\vspace{2mm}

\Q The table shows how many babies are born premature in three kinds of pregnancy.
\par\vspace{1mm}
\begin{tabular}{@{}l *{3}{c}@{}}
\toprule
Babies in the pregnancy & one & twins & triplets\\
Born premature (\%) & 7 & 57 & 93\\
\bottomrule
\end{tabular}\par\vspace{1mm}
\sub{a} Describe the pattern in the table.\lvl{M}
\alines{2}{The more babies there are in one pregnancy, the higher the percentage born premature: 7\,\% for one baby, 57\,\% for twins and 93\,\% for triplets.}
\sub{b} Suggest why babies from a multiple birth are more likely to be premature.\lvl{E}
\alines{2}{There is less space in the uterus, and the babies share the food and oxygen from the mother, so they are often born early and small.}
\sub{c} Use the table to explain why fertility treatment can be a risk for the babies.\lvl{E}
\alines{3}{Fertility treatment increases the chance of a multiple birth. More than half of twins (57\,\%) and almost all triplets are premature, and premature babies are small, with organs that are not fully developed.}

\Q A student is asked to describe IVF.\par\vspace{1mm}
\samcard{0.9\linewidth}{\flawed{In IVF the baby grows in a test tube until it is ready to be born. IVF is used when a woman's sperm ducts are blocked.}}\par\vspace{1mm}
\textbf{Sam's answer has 2 errors.} Find each one and write the correct science.\lvl{M}
\alines{4}{1. Only fertilisation and about the first two days happen in a dish; the embryo is then put into the uterus to grow (\Lref{ivf}). 2. A woman has oviducts; sperm ducts are in a man. IVF is used when the oviducts are blocked (\Lref{parts}).}
\end{document}
